{\fontsize{8.3}{10.7}\selectfont
\begin{thebibliography}{99}

\bibitem{encode2012} ENCODE Project Consortium. An integrated encyclopedia
of DNA elements in the human genome. \textit{Nature} 489, 57--74 (2012).

\bibitem{roadmap2015} Roadmap Epigenomics Consortium. Integrative analysis
of 111 reference human epigenomes. \textit{Nature} 518, 317--330 (2015).

\bibitem{bixbench2025} Mitchener, L., Laurent, J. M., Andonian, A.,
Tenmann, B., Narayanan, S., Wellawatte, G. P., White, A., Sani, L. \&
Rodriques, S. G. BixBench: a comprehensive benchmark for LLM-based agents in
computational biology. \textit{arXiv} 2503.00096 (2025).

\bibitem{compbiobench2026} Nair, S., Gunsalus, L., Orcutt-Jahns, B.,
Rossen, J., Lal, A., De Donno, C., Celik, M. H., Fletez-Brant, K., Xie, X.,
Corrada Bravo, H. \& Eraslan, G. Agentic systems are adept at solving
well-scoped, verifiable problems in computational biology. \textit{bioRxiv}
2026.04.06.716850 (2026).

\bibitem{biomysterybench2026} Anthropic. Evaluating Claude's bioinformatics
research capabilities with BioMysteryBench. Anthropic Research (2026).
\href{https://www.anthropic.com/research/Evaluating-Claude-For-Bioinformatics-With-BioMysteryBench}{anthropic.com/research/BioMysteryBench}.

\bibitem{labbench2024} Laurent, J. M., Janizek, J. D., Ruzo, M.,
Hinks, M. M., Hammerling, M. J., Narayanan, S., Ponnapati, M., White,
A. D. \& Rodriques, S. G. LAB-Bench: measuring capabilities of language
models for biology research. \textit{arXiv} 2407.10362 (2024).

\bibitem{biomnibench2026} Qu, Y., Lu, Y., Tu, X., Zhang, S., She, T.,
Shaw, A. G., Shih, J.-H., Zhao, B. et al. BiomniBench: process-level
evaluation of LLM agents for real-world biomedical research. \textit{bioRxiv}
2026.05.12.724604 (2026).

\bibitem{genebench2026} Li, J. \& Ho, A. GeneBench: assessing AI agents for
multi-stage inference problems in genomics and quantitative biology.
\textit{bioRxiv} 2026.04.22.720113 (2026).

\bibitem{diedrich2021allatac} Diedrich, J. D. et al. Profiling chromatin
accessibility in pediatric acute lymphoblastic leukemia identifies
subtype-specific chromatin landscapes and gene regulatory networks.
\textit{Leukemia} 35, 3078--3091 (2021). GEO: GSE161501.

\bibitem{barnett2023ball} Barnett, K. R. et al. Epigenomic mapping reveals
distinct B cell acute lymphoblastic leukemia chromatin architectures and
regulators. \textit{Cell Genomics} 3, 100442 (2023). GEO: GSE211631.

\bibitem{cao2020escc} Cao, W. et al. Multi-faceted epigenetic dysregulation
of gene expression promotes esophageal squamous cell carcinoma. \textit{Nature
Communications} 11, 3675 (2020). GEO: GSE149608 and GSE149609.

\bibitem{workman2025spatialbench} Workman, K., Yang, Z.,
Muralidharan, H. \& Le, H. SpatialBench: Can agents analyze real-world
spatial biology data? \textit{arXiv} 2512.21907 (2025).

\bibitem{workman2026scbench} Workman, K., Yang, Z., Muralidharan, H.,
Abdulali, A. \& Le, H. scBench: Evaluating AI agents on single-cell RNA-seq
analysis. \textit{arXiv} 2602.09063 (2026).

\bibitem{langmead2012bowtie2} Langmead, B. \& Salzberg, S. L. Fast
gapped-read alignment with Bowtie 2. \textit{Nature Methods} 9, 357--359
(2012).

\bibitem{krueger2011bismark} Krueger, F. \& Andrews, S. R. Bismark: a
flexible aligner and methylation caller for Bisulfite-Seq applications.
\textit{Bioinformatics} 27, 1571--1572 (2011).

\bibitem{buenrostro2013atacseq} Buenrostro, J. D., Giresi, P. G., Zaba,
L. C., Chang, H. Y. \& Greenleaf, W. J. Transposition of native chromatin
for fast and sensitive epigenomic profiling of open chromatin, DNA-binding
proteins and nucleosome position. \textit{Nature Methods} 10, 1213--1218
(2013).

\bibitem{zhang2008macs} Zhang, Y. et al. Model-based analysis of ChIP-Seq
(MACS). \textit{Genome Biology} 9, R137 (2008).

\end{thebibliography}
}
